%%%%%%%%%%%%%%%%%%%%%%%%%%%%%%%%%%%%%%%%%%%%%%%%%%%%%%%%%%%%%%%%%%%%%%%%%%%%
%% Author template for Management Science (mnsc) for articles with no e-companion (EC)
%% Mirko Janc, Ph.D., INFORMS, mirko.janc@informs.org
%% ver. 0.95, December 2010
%%%%%%%%%%%%%%%%%%%%%%%%%%%%%%%%%%%%%%%%%%%%%%%%%%%%%%%%%%%%%%%%%%%%%%%%%%%%
\documentclass[12pt]{article}    % <--- 12pt font
\usepackage[margin=1in]{geometry}% <--- 1 in margin
\usepackage{setspace}
\usepackage{times} % <--- times new roman
\doublespace   
%%\OneAndAHalfSpacedXII % Current default line spacing
%%\DoubleSpacedXII
%%\DoubleSpacedXI

% If hyperref is used, dvi-to-ps driver of choice must be declared as
%   an additional option to the \documentclass. For example
%\documentclass[dvips,mnsc]{informs3}      % if dvips is used
%\documentclass[dvipsone,mnsc]{informs3}   % if dvipsone is used, etc.

% Private macros here (check that there is no clash with the style)
\usepackage{color, colortbl}
\usepackage{graphicx}
%\usepackage{subfigure}
\usepackage{caption}
\usepackage{subcaption}
\usepackage{physics}
\usepackage{dsfont}
\usepackage{booktabs}
\usepackage{amsmath}


\newcommand{\gray}{\rowcolor[gray]{.90}}
% Natbib setup for author-year style
\usepackage{natbib}
 \bibpunct[, ]{(}{)}{,}{a}{}{,}%
 \def\bibfont{\small}%
 \def\bibsep{\smallskipamount}%
 \def\bibhang{24pt}%
 \def\newblock{\ }%
 \def\BIBand{and}%
\def\pipeline{prediction insight pipeline}
\newcommand{\tw}[1]{\textcolor{red}{Tong: #1}}

%% Setup of theorem styles. Outcomment only one.
%% Preferred default is the first option.
\TheoremsNumberedThrough     % Preferred (Theorem 1, Lemma 1, Theorem 2)
%\TheoremsNumberedByChapter  % (Theorem 1.1, Lema 1.1, Theorem 1.2)
\ECRepeatTheorems

%% Setup of the equation numbering system. Outcomment only one.
%% Preferred default is the first option.
\EquationsNumberedThrough    % Default: (1), (2), ...
%\EquationsNumberedBySection % (1.1), (1.2), ...

% For new submissions, leave this number blank.
% For revisions, input the manuscript number assigned by the on-line
% system along with a suffix ".Rx" where x is the revision number.

%%%%%%%%%%%%%%%%
\begin{document}

\title{On the Incorrectness and Rashomon Effect of post hoc Explanations for Business Research}

\author{Anonymous Authors}
%\date{April 2020}

% palette for colorblindness https://davidmathlogic.com/colorblind/#%23000000-%23E69F00-%2356B4E9-%23009E73-%23F0E442-%230072B2-%23D55E00-%23CC79A7



% Sample
%\KEYWORDS{deterministic inventory theory; infinite linear programming duality;
%  existence of optimal policies; semi-Markov decision process; cyclic schedule}

% Fill in data. If unknown, outcomment the field
%\KEYWORDS{post hoc, explainable, interpretable} \HISTORY{This paper was first submitted on April 12, 1922 and has been with the authors for 83 years for 65 revisions.}

\maketitle

\begin{abstract}
    Attributive post hoc explanation methods like LIME and SHAP have become very popular in machine learning. A unique problem arises when they are used in business contexts, since stakeholders may want to make decisions on post hoc explanations or draw managerial insights from feature importance measures. This is a problem because post hoc explainers have been shown to have intrinsic issues that may make their results untrustworthy. Worse yet, human evaluation studies have shown that data scientists often misuse or misunderstand post hoc explanations. In this work, we investigate how and why post hoc explanations can be misused or misunderstood. We also identify two crucial choices users make that affect the quality post hoc explanations: the choice of explainer and the choice of hyperparameters. We summarize theoretical and practical reasons for the importance of these choices and demonstrate their consequences on managerial decisions.%, and provide a framework for making them correctly. 
\end{abstract}%
%
%%%%%%%%%%%%%%%%%%%%%%%%%%%%%%%%%%%%%%%%%%%%%%%%%%%%%%%%%%%%%%%%%%%%%%

% RONILO: Idea for the structure of our critique
% 1. Identify the underlying assumptions and conditions for the use of post hoc explainers. e.g. it's hard to make an intrinsically interpretable model, small rashomonon set, already have good bbox (sunk cost fallacy...)
% 2. Presuppose the conditions for post hoc explainer ot be effective. e.g. M is accurate, data satisfies causal inference conditions. What are the effects when they are and aren't satisfied?
% 3. Examine the social relations generated by post hoc explainers. In this paper we just care about those between data scientists, model builders, and ***MANAGERS***. 
% 4. Identify the contradictions and tensions within the use of post hoc explainers. e.g. multiplicity, wrong explanations, no actionable recourse, business owners misled into making harmful decisions
% 5. Evaluate the potential for mitigation strategies involving G, M, the data, and post hoc explainers both alone and together. 
\section{Introduction}



    % Mention Marketing Science Institute "consistently identified attribution as a topmost research priority over the last few years" (from SHAP meets uniform paper)

% =========================== Tong ===========================
In the past decade, the volume and complexity of data has significantly increased, leading to the widespread adoption of machine learning for business analysis.


However, business researchers prioritize gaining managerial insights rather than making accurate predictions as in typical machine learning.


Therefore, a new class of methods from machine learning naturally caught the attention of business researchers: \emph{post hoc} explainers.

The business research community readily adopts this new type of method to understand their black-box models and obtain managerial insights. 
    \begin{figure}[t!]
        \centering
        \includegraphics[width = \textwidth]{Figures/pipeline.png}
        \caption{Prediction Insight Pipeline: the pipeline of using post hoc explainers to obtain managerial insights in business research.}
        \label{fig:pipeline}
    \end{figure}

Although post hoc explanations seem like a magical tool via which researchers and practitioners can enjoy both high predictive accuracy and explainability at the same time, an alarm has gone off in the machine learning community.

Despite the heated debate about the usage of post hoc explainers, the discussion in the machine learning community is mainly focused on one leg of the pipeline in Figure \ref{fig:pipeline}  (``explaining a black-box''). 

This paper explores how post hoc explanations could lead to incorrectness in two specific types of descriptive and prescriptive managerial insights. 
    
\section{Related Work}
    \label{sec:related_work}
    This section introduces popular post hoc explainers and then give sa review of prior literature discussing the weaknesses of post hoc explainers. 

    \subsection{Background on LIME, SHAP and DiCE}
    There are two main families of post hoc explainers: \textbf{attributive} and \textbf{counterfactual}. 
    
    \paragraph{Attributive Explanations} Given a black box model $\mathcal{M}$ taking an input $x=(x_1,...,x_p)$ with $p$ features, attributive explainers explain $\mathcal{M}(x)$ by assigning a feature importance score to each of the $p$ features.


 \paragraph{Counterfactual Explanations}   Counterfactual explainers explain $\mathcal{M}(x)$ by identifying minimal changes to the feature values $x_i$ that would result in an output different from $\mathcal{M}(x)$.

    \subsection{Weaknesses of LIME, SHAP, and DiCE}
        \label{sec:weaknesses}
        As aforementioned, there are three common weaknesses that post hoc explainers possess. 
    

        % INSTABILITY
        % ADD TRANSITION SENTENCE HERE
    Post hoc explainers make a trade-off between stability and local accuracy \citep{alvarez18robustness, laugel18defining}. 


    % consistency
    In addition to having the intrinsic issue of instability, explainers are also sensitive to externalities decided by the user. 
    % CORRECTNESS AND EFFECTIVENESS
    % contrast by saying e.g. kumar was only about stage 1
    Potentially worse than instability and inconsistency is the tendency for post hoc explainers' explanations to be incorrect and/or ineffectual. \citep{Krishna2022TheDP,kumar20shap,weerts19human,dieber20assessing,yacoby22judges}. 
    


\section{Simulated Dataset}
\label{sec:dataset}
In this section, we will discuss a simulated loan approval dataset that we will use to benchmark how post-hoc explainers work in a hypothetical business scenario.
 
 \paragraph{Data Generator $\mathcal{G}$}

Table \ref{tab:simulated} summarizes the meanings of all the features used by $\mathcal{G}$ and how they were generated. 
 
 

 \begin{table}[h]
\centering
\resizebox{\textwidth}{!}{%
\begin{tabular}{|l|l|l|l|}
\hline
\textbf{Feature} & \textbf{Explanation}                           & \textbf{Formula}       & \textbf{Comment}                             \\ \hline
%\rowcolor[gray]{.9}
Age              & Applicant Age                                   & Uniform(18,100)        &                                              \\ \hline
%\rowcolor[gray]{.9}
Sex              & Applicant Sex                                   & Bernoulli(.5)          &                                              \\ \hline
Married          & Applicant Marital Status                        & Bernoulli(.5)          &                                              \\ \hline
Social Support   & 0 = no support, 1 = perfectly supported        & Uniform(0,1)           & Not observable                               \\ \hline
Cosigner         & Indicator of presence of loan cosigner         & Bernoulli(.5)          & Depends on social support and marital status \\ \hline
%Economy          & Health of the economy in the customer's region & Uniform(0,1)           &                                %              \\ \hline
Income           & Applicant Income                                & Normal(40,000, 30,000) & %Colinear with Economy 
\\ \hline
%\rowcolor[gray]{.9}
Utilization      & Proportion of total credit utilized            & Uniform(0,1)           & %Colinear with Economy 
\\ \hline
Delinquencies    & Number months since of last deliquent payment                   & Uniform(0,25)          &                                              \\ \hline
%\rowcolor[gray]{.9}
Inquiries         & Number of credit inquiries                     & Uniform(0,6)           &                                              \\ \hline
%\rowcolor[gray]{.9}
Credit Score     & Credit score based on FICO scoring             & Given below            & Depends on inquiries, delinquencies, and utilization.                                             \\ \hline
\end{tabular}%
}
\caption{Description of features for simulated loan data. }%Note that for demonstration purposes, Section \ref{sec:pipeline_demo} experiments use all features but the unobserved social support variable.}
\label{tab:simulated}
\end{table}



All experiments assume there is an underlying true model $\mathcal{G}$ that classifies whether individuals are approved for a loan. 



\section{Incorrectness of Post Hoc Explanations}
    In this section, we investigate the correctness of managerial insights derived from post hoc explanations. 


            \subsection{Unfaithful Descriptive Insights}
                    \label{sec:unfaithful_descriptive}
To assess how well LIME and SHAP can provide managers with useful descriptive insights, we investigate how faithful LIME and SHAP explanations are to $\mathcal{G}$. 

                     
Throughout this paper, we assess the correctness of feature importance rankings of LIME and SHAP using the notion of Spearman Rank Correlation \citep{MyersJeromeL2010Rdas}.
                     
 In Section \ref{sec:relative_importance}, we will compute the Spearman Rank Correlations between the ground truth feature ranking in $\mathcal{G}$ (formula \ref{eqn:gt}) and those given by LIME and SHAP for 100 randomly sampled loan applicants.



                    \subsubsection{Identifying Relative Importance}
                        \label{sec:relative_importance}
                        We can tell the relative importance of features from LIME by comparing the magnitudes of its surrogate model's coefficients. 


                    % GIVE DEFINITION OF RECALL
                    \subsubsection{Recall of the Most Important Features}
                        The previous section explores the abilities of LIME and SHAP to recover correct importance information about all features used by $\mathcal{G}$, but in practice  users may only care about, say, the top 5 most important features. 

                        Finding the top $K$ most important features to a model is one of the most common use cases for LIME and SHAP, so the fact that they cannot successfully do it reliably should urge business researchers to use LIME and SHAP with caution for this purpose.
                        
          \begin{figure}[ht!]
            \centering
            \includegraphics[width=.5\textwidth]{Figures/lime_shap_topk.png}
            \caption{Average recall of LIME and SHAP when attempting to capture the top K most important features to $\mathcal{G}$ when explaining $\mathcal{M}$, for varied $K=1,2,3,...\#\text{Features}$. Vertical lines indicate the 95\% confidence interval of the recall for given $K$.}
            \label{fig:avg_recall}
        \end{figure}  


    \subsection{Unfaithful Prescriptive Insight}
        \label{sec:actionable_recourse}
        Since section \ref{sec:unfaithful_descriptive} showed that LIME and SHAP may fail both to correctly indicate what features are most important to loan approval and how the features affect loan approval probability, this section will show that they may additionally fail to be useful in advising how to change features to obtain more desirable outcomes.
        
        We test two ways to obtain post hoc prescriptive insights.

            
            In our simulation experiment, we suppose that applicants use an MLP $\mathcal{M}$ trained on our data set to be able to anticipate banks' loan approval decisions. 

            
            To show that using LIME and SHAP for actionable recourse may fail, we perturb the top $K$ most important features of 50 denied applicants in the simulated loan approval data set (as seen by LIME and SHAP) and track the percentage of true counterfactuals that LIME and SHAP can produce.


    

\begin{figure}[!htb]
\minipage{0.32\textwidth}
  \includegraphics[width=\linewidth]{Figures/lime_recourse_k.png}
  \label{fig:lime_recourse_k}
  \caption{A really Awesome Image}\label{fig:awesome_image1}
\endminipage\hfill
\minipage{0.32\textwidth}
  \includegraphics[width=\linewidth]{Figures/shap_recourse_k.png}
  \label{fig:lime_recourse_k}
  \caption{A really Awesome Image}\label{fig:awesome_image2}
\endminipage\hfill
\minipage{0.32\textwidth}%
  \includegraphics[width=\linewidth]{Final_Figures/fig8a.png}
  \label{fig:dice_recourse}
  \caption{A really Awesome Image}\label{fig:awesome_image3}
\endminipage
\end{figure}
        
 %\tw{Ronilo experiment: allow changing top K features and report the successful recourse for LIME and SHAP}           
            Figures 3 and 4 plot the perturbation extent versus the proportion of times previously rejected applicants get approved by both $\mathcal{M}$ and $\mathcal{G}$ for LIME and SHAP, respectively.

            Although post hoc explainers may seem to present a way to develop interventions or provide explanations to customers or stakeholders via counterfactual generation, doing so does not seem to be straightforward. 



\section{Rashomon Effect of Post Hoc Explanations}
        A Rashomon effect of post hoc explainers refers to the phenomenon in which post hoc explainers provide a set of explanations that is inconsistent such that one explainer contradicts itself, two explainers contradict each other, or both.
        
        Due to the Rashomon effect, researchers cannot just choose one explainer or one explanation. 

       \subsection{Inconsistencies Between Explainers}
        \label{sec:conflict_between}
            %To motivate the use of multiple post hoc explainers, we show that they may explain the same outcome in different ways with differing degrees of accuracy. 
            Conflicting explanations between explainers can lead to a conundrum for researchers. 
    
            To justify these claims, we demonstrate the inconsistency of LIME and SHAP on the simulated loan data set.  

            % REWORK NEXT TO PARAGRAPHS
            These results suggest once more that post hoc explainers should be used with caution.

            \begin{figure}[ht!]
                \centering
                \includegraphics[width=.5\textwidth]{Figures/lime_shap_spearman_together.png}
                \caption{Violin plot of Spearman Rank Correlations between LIME and SHAP explanations of 100 points}
                \label{fig:lime_shap_spearman_together}
            \end{figure}


    \subsection{Inconsistency from the Same Explainer}
            \label{sec:conflict_within}
             Both LIME and SHAP are designed such that there may be variability in explanations of the same individual between runs. 

        

    \section{Attempted "Mitigation" Strategies}
        \subsection{Attempted ``Mitigation'' Strategies for Incorrectness}
        \label{sec:mitigation} %\textcolor{red}{move this to the end of section 4. this is the mitigation strategy for incorrectness}
            We propose some mitigation strategies based on tuning $\mathcal{M}$ and trying to take advantage of the inconsistency laid out in the previous section.


            \subsubsection{Does High-Accuracy of $\mathcal{M}$ Imply High-Quality Explanations?}
            \label{sec:m_fidelity}
                To answer this question, we used two strategies to control the accuracy of the model.
            
                \begin{figure}[ht!]
                \centering
                    \begin{subfigure}{.48\textwidth}
                          \centering
                          \includegraphics[width=\linewidth]{Final_Figures/fig14a.png}%{Figures/feng_fig/acc_increase_iter_1.png}
                          %\caption{A subfigure}
                          \label{fig:acc_increase_metric1_strategy1}
                    \end{subfigure}%
                    \hfill
                    \begin{subfigure}{.48\textwidth}
                          \centering
                          \includegraphics[width=\linewidth]{Final_Figures/fig14b.png}%{Figures/feng_fig/acc_increase_size_1.png}
                          %\caption{A subfigure}
                          \label{fig:acc_increase_metric1_strategy2}
                    \end{subfigure}
                    \hfill
                \caption{Line plots of 100 Spearman rank correlations of feature importance magnitude ranking given by LIME and SHAP when increasing $\mathcal{M}'s$ accuracy by changing model parameter strategy (left) and changing training size (right) strategy.}
                \label{fig:acc_increase_metric1}
                \end{figure}
        \subsection{Attempted ``Mitigation'' Strategy for Inconsistency} 
        \label{sec:mitigation}
            Section \ref{sec:conflict_within} demonstrated the multiplicity of explanations due to random seeds and hyperparameters, so it is natural to wonder if multiple explanations can be aggregated to provide an explanation that tends to be the most accurate. 
            
            \subsubsection{Averaging Explanations of One Explainer}
                We check if taking the mean of feature importance values from LIME and SHAP taken over combinations of random seeds and hyperparameters reveals a useful aggregation strategy. 

            
            \subsubsection{Averaging Explanations Across Explainers}
                Since \citep{kommiya21towards} demonstrated that feature importance rankings differed significantly across different explainers, we also seek a useful aggregation strategy for explanations from LIME and SHAP. 

                Figure \ref{fig:lime_shap_Spearman_together} shows the results of this experiment, and it is not clear that aggregating LIME and SHAP is beneficial. 


\begin{figure}[!htb]
\minipage{0.32\textwidth}
  \includegraphics[width=\linewidth]{Final_Figures/fig22a.png}
  %\label{fig:avering_within_lime_metric1}
  \caption{Violin plots of 100 Spearman rank correlations of LIME feature importance magnitude rankings when averaging different numbers of model results}\label{fig:avering_within_lime_metric1}
\endminipage\hfill
\minipage{0.32\textwidth}
  \includegraphics[width=\linewidth]{Final_Figures/fig22b.png}
  %\label{fig:lime_recourse_k}
  \caption{Violin plots of 100 Spearman rank correlations of LIME feature importance magnitude rankings when averaging different numbers of model results}
  \label{fig:avering_within_shap_metric1}
\endminipage\hfill
\minipage{0.32\textwidth}%
  \includegraphics[width=\linewidth]{Final_Figures/fig20.png}
  %\label{fig:dice_recourse}
  \caption{Violin plots of Spearman Rank Correlations of feature importance magnitude ranking by LIME, SHAP and LIME+SHAP}
  \label{fig:lime_shap_Spearman_together}
\endminipage
\end{figure}

                

  

\section{Conclusion}
    Model-agnostic post hoc explainers certainly have legitimate use cases, but due to the nature of real business data sets and stakes associated with business problems, we urge the business community to use them with caution and only when necessary. 

    Throughout this paper we have shown various consequences of the misuse of post hoc explainers that could lead to losses for both businesses and their customers. 
        



\bibliographystyle{informs2014}
\bibliography{refs}

\newpage

% dimensions finds only 206 but indexes preprints
% scopus finds over 400 but doesn't index preprints
%% preprint not necessary
% need to ensure that the papers are actually explaining a black box
% e.g. add

%%%%%%%%%%%%%%%%%
\end{document}
%%%%%%%%%%%%%%%%%

